\section{Results}
\label{sec:results}

\subsection{Metrics and Decision Rule}
\label{sec:metrics}

\paragraph{Removal.}
We report the removal rate at $\alpha = 0.05$ under the generating key, with a $95\%$ prompt-clustered bootstrap interval and a document-level Wilson interval beside it, and the corpus-level true-positive rate at a $1\%$ false-positive rate from the parametric null. The first is the per-document question, the second is the detector-tuning question, and they can disagree: a chain can push most documents below the per-document threshold while leaving a detectable residual in the pooled distribution. Reporting both makes the disagreement visible.

\paragraph{Preservation.}
Two signals form the selecting composite. BERTScore-F1 measures embedding-level similarity to the source \citep{bertscore}. COMET measures adequacy and fluency, with the source document on the source side, the attacked document as the hypothesis, and the watermarked source as the reference \citep{comet}. Both are normalised to $[0,1]$, and $S$ is their unweighted mean. BLEURT provides a second learned judgement \citep{bleurt} and MAUVE measures distributional closeness between the attacked and unattacked sets \citep{mauve}; both are reported on the same documents and do not enter $S$ or $\tau$, so adding BLEURT at the larger tier cannot move which chain is selected. We also report the length ratio and the achieved edit rate, because a tactic that shortens the document can inflate similarity scores while dropping content. Per-metric values are always reported so a reader can reweight them, and the fraction of documents below $\tau$ is reported for every recommended chain.

\paragraph{Threshold.}
The semantic floor is $\tau = 0.80$, fixed before the search and revised at most once after the pilot run, with any revision recorded as a dated deviation in the preregistration. Results at $\tau - 0.05$ and $\tau + 0.05$ are reported in Appendix~\ref{app:perdomain} so that the sensitivity of the recommendation to the floor is visible. The attack prompts are in Appendix~\ref{app:prompts}.

\subsection{Single Tactics}

Table~\ref{tab:singles} reports every tactic applied once. The control row is the reference point: its removal rate should match the unattacked baseline.

\begin{table}[t]
\centering
\small
\caption{Single tactics at their default intensity. Higher removal is better; higher TPR at $1\%$ FPR means the detector still finds the mark. Preservation columns are higher-is-better except length ratio, which is closest to $1$ at $1.0$.}
\label{tab:singles}
\setlength{\tabcolsep}{4pt}
\begin{tabular}{@{}>{\raggedright\arraybackslash}p{4.0cm}ccccc@{}}
\toprule
Tactic & Removal $\uparrow$ & Residual TPR@1\%FPR $\downarrow$ & BERTScore $\uparrow$ & COMET $\uparrow$ & Length ratio \\
\midrule
Instruction paraphrase   & \ph{run} & \ph{run} & \ph{run} & \ph{run} & \ph{run} \\
Discourse paraphrase     & \ph{run} & \ph{run} & \ph{run} & \ph{run} & \ph{run} \\
Masked-LM infilling      & \ph{run} & \ph{run} & \ph{run} & \ph{run} & \ph{run} \\
Style normalisation      & \ph{run} & \ph{run} & \ph{run} & \ph{run} & \ph{run} \\
Document rewrite         & \ph{run} & \ph{run} & \ph{run} & \ph{run} & \ph{run} \\
Single-pivot back-translation & \ph{run} & \ph{run} & \ph{run} & \ph{run} & \ph{run} \\
Multi-hop back-translation & \ph{run} & \ph{run} & \ph{run} & \ph{run} & \ph{run} \\
Cross-lingual generation & \ph{run} & \ph{run} & \ph{run} & \ph{run} & \ph{run} \\
Synonym substitution     & \ph{run} & \ph{run} & \ph{run} & \ph{run} & \ph{run} \\
Random word edits        & \ph{run} & \ph{run} & \ph{run} & \ph{run} & \ph{run} \\
Typo and noise injection & \ph{run} & \ph{run} & \ph{run} & \ph{run} & \ph{run} \\
Non-origin regeneration  & \ph{run} & \ph{run} & \ph{run} & \ph{run} & \ph{run} \\
\midrule
Unicode scrub (control)  & \ph{run} & \ph{run} & \ph{run} & \ph{run} & \ph{run} \\
\bottomrule
\end{tabular}
\end{table}

\subsection{Pairs and Triples}

Tables~\ref{tab:pairs} and~\ref{tab:triples} report the best chains at each depth on the held-out split, ordered by removal rate among chains that satisfy the semantic floor. The full ranked list is released with the results archive.

\begin{table}[t]
\centering
\small
\caption{Best ordered pairs. Ranks are by removal rate among admissible chains.}
\label{tab:pairs}
\begin{tabular}{@{}clcccc@{}}
\toprule
Rank & Chain & Removal $\uparrow$ & Residual TPR@1\%FPR $\downarrow$ & COMET $\uparrow$ & Length ratio \\
\midrule
1 & \ph{chain} & \ph{run} & \ph{run} & \ph{run} & \ph{run} \\
2 & \ph{chain} & \ph{run} & \ph{run} & \ph{run} & \ph{run} \\
3 & \ph{chain} & \ph{run} & \ph{run} & \ph{run} & \ph{run} \\
4 & \ph{chain} & \ph{run} & \ph{run} & \ph{run} & \ph{run} \\
5 & \ph{chain} & \ph{run} & \ph{run} & \ph{run} & \ph{run} \\
\bottomrule
\end{tabular}
\end{table}

\begin{table}[t]
\centering
\small
\caption{Best ordered triples. The third step is compared against the pair it extends.}
\label{tab:triples}
\begin{tabular}{@{}clcccc@{}}
\toprule
Rank & Chain & Removal $\uparrow$ & Residual TPR@1\%FPR $\downarrow$ & COMET $\uparrow$ & Length ratio \\
\midrule
1 & \ph{chain} & \ph{run} & \ph{run} & \ph{run} & \ph{run} \\
2 & \ph{chain} & \ph{run} & \ph{run} & \ph{run} & \ph{run} \\
3 & \ph{chain} & \ph{run} & \ph{run} & \ph{run} & \ph{run} \\
4 & \ph{chain} & \ph{run} & \ph{run} & \ph{run} & \ph{run} \\
5 & \ph{chain} & \ph{run} & \ph{run} & \ph{run} & \ph{run} \\
\bottomrule
\end{tabular}
\end{table}

\subsection{Frontier}

Figure~\ref{fig:pareto} plots removal against preservation for every evaluated chain. The frontier is the set of chains not dominated on both axes. The recommendation in \S\ref{sec:discussion} is the frontier point with the highest removal among chains at or above $\tau$.

\begin{figure}[t]
\centering
\fbox{\parbox[c][0.28\textheight][c]{0.9\linewidth}{\centering \ph{Pareto plot: removal rate on the vertical axis, the selecting composite $S$ (BERTScore and COMET) on the horizontal axis; single tactics, pairs, and triples in separate series; the semantic floor drawn as a vertical line.}}}
\caption{Removal against preservation for all evaluated chains.}
\label{fig:pareto}
\end{figure}

\subsection{Second-Attacker Ablation}

The top five chains, the best single tactic, and the control are re-run end to end with Claude Sonnet on the held-out split. Table~\ref{tab:ablation} reports removal under both attackers and the change. The rank correlation across the seven rows is $\rho = \ph{value}$, and the largest absolute removal change is \ph{value}. \ph{One sentence on whether the ranking transfers and what the exception is.}

\subsection{Per-Domain Results}

Table~\ref{tab:perdomain} repeats the recommended chain per domain. Domain effects are expected: news writing constrains entities and numbers, question answering constrains facts, and creative writing leaves more room for paraphrase.

\begin{table}[t]
\centering
\small
\caption{Recommended chain per domain at the same floor $\tau$.}
\label{tab:perdomain}
\begin{tabular}{@{}llcc@{}}
\toprule
Domain & Recommended chain & Removal $\uparrow$ & COMET $\uparrow$ \\
\midrule
News writing              & \ph{chain} & \ph{run} & \ph{run} \\
Long-form QA              & \ph{chain} & \ph{run} & \ph{run} \\
Creative writing          & \ph{chain} & \ph{run} & \ph{run} \\
MarkMyWords tasks         & \ph{chain} & \ph{run} & \ph{run} \\
\bottomrule
\end{tabular}
\end{table}
